%% kapa: exam=1 round=37 year=2026 period=2 form=A subject=appraisal_law subject_ko=감정평가관계법규 seq=6 number=6
%% source: Q-Net 감정평가사 기출문제 (cst00309) artlSeq=5258924 file=raw/1008/5258924/2026년 제37회 감정평가사 1차 2교시.pdf
%% conversion: kordoc 4.24.51 -> tools/gen1.py (auto); review=proofread-sonnet; flags=none
\begin{question}{6}
국토의 계획 및 이용에 관한 법령상 개발행위에 따른 기반시설의 설치에 관한 설명으로 옳지 않은 것은?
\begin{choices}
  \item 개발밀도관리구역의 범위를 변경하려면 해당 지방자치단체에 설치된 지방도시계획위원회의 심의를 거쳐야 한다.
  \item 개발밀도관리구역에서는 용도지역에 따른 건폐율 또는 용적률을 완화하여 적용한다.
  \item 기반시설부담구역에서 200제곱미터를 초과하는 숙박시설의 증축 행위는 기반시설설치 비용의 부과대상이 된다.
  \item 국토교통부장관은 매년 1월 1일을 기준으로 한 기반시설 표준시설비용을 매년 6월 10일까지 고시하여야 한다.
  \item 기반시설부담구역에서 기반시설설치비용의 부과대상인 건축행위를 위탁한 경우 그 위탁을 한 자는 기반시설설치비용을 내야 한다.
\end{choices}
\end{question}
\answer{2}
